% ==============================================================================
% CHAPTER 7 RESULTS - PLACEHOLDER TEMPLATE
% ==============================================================================
% This template provides the STRUCTURE and INTERPRETATION NARRATIVE for
% Chapter 7 results, fully writeable NOW without waiting for training completion.
% Numbers will be filled when jobs complete.
% 
% INSERT THIS INTO: LaTeX_Thesis/chapters/07_results.tex
% ==============================================================================

\chapter{Results}
\label{ch:results}

This chapter presents the experimental results from stages 3 and 4 of the KVP extraction pipeline. Stage 3 establishes a text-only generative baseline (Mistral-7B); Stage 4 introduces layout-aware discriminative extraction (LayoutLMv3) to validate our hypothesis that spatial grounding improves extraction accuracy, particularly on sparse document layouts.

\section{Stage 3: Mistral-7B Generative Baseline}
\label{sec:results_stage3}

Stage 3 implements the official KVP10k baseline (Mistral-7B-Instruct, fine-tuned with LoRA on 4,985 documents) and evaluates its performance on the held-out test set (581 unique document pages with valid KVP annotations).

\subsection{Overall Performance}
\label{sec:stage3_overall}

Table~\ref{tab:stage3_overall} summarizes end-to-end performance under two evaluation protocols: the strict Stage-0 protocol (NED$<0.2$, IoU$>0.3$) and the default threshold protocol (NED$<0.5$, IoU$>0.5$).

\begin{table}[ht]
\centering
\caption{Stage 3 (Mistral-7B) Overall Performance}
\label{tab:stage3_overall}
\begin{tabular}{@{}lccc@{}}
\toprule
\textbf{Protocol} & \textbf{Precision} & \textbf{Recall} & \textbf{F1} \\
\midrule
Stage-0 (NED$<0.2$, IoU$>0.3$) & \texttt{[TBD]} & \texttt{[TBD]} & \texttt{[TBD]} \\
Default (NED$<0.5$, IoU$>0.5$) & \texttt{[TBD]} & \texttt{[TBD]} & \texttt{[TBD]} \\
\bottomrule
\end{tabular}
\end{table}

\textit{Interpretation}: [Will be computed from training logs] These results represent the performance of a pure text-based generative model without spatial reasoning. The generative paradigm allows open-vocabulary KVP extraction but introduces hallucination risk, characterized by plausible-sounding but unsupported predictions.

\subsection{Performance by Layout Density}
\label{sec:stage3_by_cluster}

A key finding from Stage 3 is that performance varies dramatically with document layout density. Test documents were grouped into two clusters based on spatial density of text fields (see Section~\ref{sec:clustering} for cluster definitions). Table~\ref{tab:stage3_cluster_performance} reports performance breakdown by cluster.

\begin{table}[ht]
\centering
\caption{Stage 3 Performance by Layout Cluster (Stage-0 Protocol)}
\label{tab:stage3_cluster_performance}
\begin{tabular}{@{}lccccc@{}}
\toprule
\textbf{Cluster} & \textbf{Type} & \textbf{N} & \textbf{F1} & \textbf{Correct} & \textbf{Hallucinate Rate} \\
\midrule
0 & Sparse & \texttt{[TBD]} & \texttt{[TBD]} & \texttt{[TBD]} & \texttt{[TBD]} \\
1 & Dense & \texttt{[TBD]} & \texttt{[TBD]} & \texttt{[TBD]} & \texttt{[TBD]} \\
\bottomrule
\end{tabular}
\end{table}

\textit{Interpretation}: [Will be computed] We expect large performance variation between clusters, with sparse layouts showing significantly lower correctness and higher hallucination rates. This asymmetry motivates Stage 4: layout-aware spatial models should leverage clustering information to reduce hallucination on sparse documents while maintaining performance on dense layouts.

\subsection{Error Analysis: Hallucination vs.\ Omission}
\label{sec:stage3_error_analysis}

Stage 3 errors fall into three categories:

\begin{enumerate}
    \item \textbf{Correct}: Predicted KVP matches ground truth (text and location both accurate).
    
    \item \textbf{Hallucinated}: Predicted KVP has no correspondence in ground truth, or corresponds but falls outside spatial or textual tolerance thresholds.
    
    \item \textbf{Missed}: Ground-truth KVP was not predicted by the model.
\end{enumerate}

Table~\ref{tab:stage3_error_breakdown} reports the distribution of errors by type and cluster.

\begin{table}[ht]
\centering
\caption{Stage 3 Error Breakdown by Cluster}
\label{tab:stage3_error_breakdown}
\begin{tabular}{@{}lccccc@{}}
\toprule
\textbf{Cluster} & \textbf{Sample Count} & \textbf{\% Correct} & \textbf{\% Hallucinate} & \textbf{\% Missed} \\
\midrule
0 (Sparse) & \texttt{[TBD]} & \texttt{[TBD]}\% & \texttt{[TBD]}\% & \texttt{[TBD]}\% \\
1 (Dense) & \texttt{[TBD]} & \texttt{[TBD]}\% & \texttt{[TBD]}\% & \texttt{[TBD]}\% \\
\midrule
Overall & \texttt{[TBD]} & \texttt{[TBD]}\% & \texttt{[TBD]}\% & \texttt{[TBD]}\% \\
\bottomrule
\end{tabular}
\end{table}

\textit{Interpretation}: [Will be computed] We hypothesize that sparse layouts suffer from high hallucination (Mistral over-generates spurious predictions) while dense layouts suffer more from omissions (many correct fields are present but not extracted). Text-only models lack spatial cues to distinguish spurious vs.\ genuine field locations.

\section{Stage 4: LayoutLMv3 Discriminative Models}
\label{sec:results_stage4}

Stage 4 replaces the generative pipeline with a discriminative architecture combining LayoutLMv3 encoding with token-level entity classification (Stage 4a) and learned key-value linking (Stage 4b). This section presents results from both stages and the $\lambda$ sweep over relation weight.

\subsection{Stage 4a: Entity Classification Only}
\label{sec:stage4a_results}

Stage 4a predicts token-level entity labels (Key, Value, Other) without explicit relation learning. Table~\ref{tab:stage4a_overall} reports performance.

\begin{table}[ht]
\centering
\caption{Stage 4a (Entity-Only) Overall Performance}
\label{tab:stage4a_overall}
\begin{tabular}{@{}lccc@{}}
\toprule
\textbf{Protocol} & \textbf{Precision} & \textbf{Recall} & \textbf{F1} \\
\midrule
Stage-0 (NED$<0.2$, IoU$>0.3$) & \texttt{[TBD]} & \texttt{[TBD]} & \texttt{[TBD]} \\
Default (NED$<0.5$, IoU$>0.5$) & \texttt{[TBD]} & \texttt{[TBD]} & \texttt{[TBD]} \\
\bottomrule
\end{tabular}
\end{table}

\textit{Interpretation}: [Will be computed] Stage 4a should show improvement over Stage 3 overall (spatial grounding helps), and especially large improvements on Cluster 0 (sparse layouts). If Stage 4a+spatial is already strong, the linking problem is not the bottleneck on this task.

\subsection{Stage 4b: Entity + Relation Linking}
\label{sec:stage4b_results}

Stage 4b adds a biaffine scorer to predict key-value relationships, trained with weight $\lambda$ on the relation loss. Table~\ref{tab:stage4b_lambda_sweep} presents the $\lambda$ sweep results.

\begin{table}[ht]
\centering
\caption{Stage 4b Entity+Relation Performance (\texorpdfstring{$\lambda$}{lambda} Sweep)}
\label{tab:stage4b_lambda_sweep}
\begin{tabular}{@{}lccccc@{}}
\toprule
\textbf{Lambda} & \textbf{Precision} & \textbf{Recall} & \textbf{Entity F1} & \textbf{Relation F1} & \textbf{Overall F1} \\
\midrule
0.5 (Entity-weighted) & \texttt{[TBD]} & \texttt{[TBD]} & \texttt{[TBD]} & \texttt{[TBD]} & \texttt{[TBD]} \\
1.0 (Balanced) & \texttt{[TBD]} & \texttt{[TBD]} & \texttt{[TBD]} & \texttt{[TBD]} & \texttt{[TBD]} \\
2.0 (Relation-weighted) & \texttt{[TBD]} & \texttt{[TBD]} & \texttt{[TBD]} & \texttt{[TBD]} & \texttt{[TBD]} \\
\bottomrule
\end{tabular}
\end{table}

\textit{Interpretation}: [Will be computed] If $\lambda=1.0$ performs best, entity and relation learning are equally important. If $\lambda=0.5$ is best, the task is entity-dominated (relations are learned easily). If $\lambda=2.0$ is best, relation learning is the critical bottleneck.

\subsection{Performance by Cluster (Stage 4a and 4b)}
\label{sec:stage4_by_cluster}

Table~\ref{tab:stage4_cluster_performance} shows how discriminative models perform across layout clusters, comparing Stage 4a (entity only) and Stage 4b (entity + relation) at best $\lambda$.

\begin{table}[ht]
\centering
\caption{Stage 4 Performance by Layout Cluster}
\label{tab:stage4_cluster_performance}
\begin{tabular}{@{}lcccccc@{}}
\toprule
\textbf{Stage} & \textbf{Cluster} & \textbf{N} & \textbf{F1} & \textbf{\% Correct} & \textbf{\% Hallucinate} & \textbf{\% Missed} \\
\midrule
4a & 0 (Sparse) & \texttt{[TBD]} & \texttt{[TBD]} & \texttt{[TBD]} & \texttt{[TBD]} & \texttt{[TBD]} \\
4a & 1 (Dense) & \texttt{[TBD]} & \texttt{[TBD]} & \texttt{[TBD]} & \texttt{[TBD]} & \texttt{[TBD]} \\
4b & 0 (Sparse) & \texttt{[TBD]} & \texttt{[TBD]} & \texttt{[TBD]} & \texttt{[TBD]} & \texttt{[TBD]} \\
4b & 1 (Dense) & \texttt{[TBD]} & \texttt{[TBD]} & \texttt{[TBD]} & \texttt{[TBD]} & \texttt{[TBD]} \\
\bottomrule
\end{tabular}
\end{table}

\textit{Interpretation}: [Will be computed] We expect that spatial grounding eliminates most hallucinations on Cluster 0. Whether Stage 4b (with linking) further improves over 4a will reveal whether the linking problem is learnable.

\section{Stage 3 vs.\ Stage 4 Comparison}
\label{sec:stage3_vs_stage4}

Table~\ref{tab:stage3_vs_stage4_comparison} compares the generative baseline (Stage 3, Mistral-7B, text-only) against the best discriminative model (Stage 4b with optimal $\lambda$, LayoutLMv3, layout-aware).

\begin{table}[ht]
\centering
\caption{Overall Comparison: Stage 3 (Text-Only Generative) vs.\ Stage 4 (Layout-Aware Discriminative)}
\label{tab:stage3_vs_stage4_comparison}
\begin{tabular}{@{}lcccc@{}}
\toprule
\textbf{Model} & \textbf{Overall F1} & \textbf{Cluster 0 (Sparse)} & \textbf{Cluster 1 (Dense)} & \textbf{Improvement (overall)} \\
\midrule
Stage 3 (Mistral) & \texttt{[TBD]} & \texttt{[TBD]} & \texttt{[TBD]} & --- \\
Stage 4b (LayoutLMv3) & \texttt{[TBD]} & \texttt{[TBD]} & \texttt{[TBD]} & \texttt{[TBD]}\% \\
\bottomrule
\end{tabular}
\end{table}

\textit{Interpretation}: [Will be computed] Stage 4 should show substantial improvement on Cluster 0 (sparse), driven by elimination of spatial hallucination. Overall improvement will depend on whether spatial grounding benefits Cluster 1 (dense layouts).

\section{Ablation Study: Architecture Components}
\label{sec:ablation_results}

Section~\ref{sec:stage4_ablation} outlined three ablations:

\begin{enumerate}
    \item \textbf{4a vs 4b}: Does entity-only learning suffice, or does relation linking help?
    \item \textbf{$\lambda$ sweep (0.5, 1.0, 2.0)}: How should entity and relation losses be balanced?
    \item \textbf{Stage 3 vs 4}: How much of the improvement comes from spatial vs.\ architectural changes?
\end{enumerate}

\subsection{Ablation 1: Entity-Only vs.\ Entity+Relation}
\label{sec:ablation_entity_vs_relation}

Figure~\ref{fig:ablation_4a_vs_4b} visualizes the contribution of relation learning. Comparing Stage 4a (no linking) against Stage 4b (biaffine linking) reveals whether the linking component is a bottleneck.

\textit{Interpretation}: [Will be computed and visualized] If 4b substantially improves over 4a (e.g., 5\% F1 gain), learned linking is valuable. If 4b $\approx$ 4a, entity classification is the limiting factor.

\subsection{Ablation 2: Loss Weight Sweep}
\label{sec:ablation_lambda_sweep}

Figure~\ref{fig:lambda_sweep} shows F1 across $\lambda \in \{0.5, 1.0, 2.0\}$ for Stage 4b. This sweep determines whether entity or relation learning dominates.

\textit{Interpretation}: [Will be computed] If $\lambda=1.0$ is best, both objectives matter equally. If $\lambda=0.5$ is best, over-emphasizing relations hurts performance. If $\lambda=2.0$ is best, relation learning is critical.

\subsection{Ablation 3: Spatial Grounding Contribution}
\label{sec:ablation_spatial}

Comparing Stage 3 (text-only) vs.\ Stage 4 (text+layout) directly isolates the contribution of spatial information. However, confounding factors (generative vs.\ discriminative, architecture size) cannot be controlled, so this comparison is observational rather than causal.

\textit{Interpretation}: [Will be computed] If Stage 4 >> Stage 3, spatial grounding is important. The per-cluster breakdown reveals whether spatial benefits are distributed or concentrated on particular layout types.

\section{Robustness Analysis}
\label{sec:robustness}

\textit{Note}: Robustness analysis is deferred to Phase E (future work) and is not included in the current evaluation. Future work should include adversarial robustness, domain shift analysis, and sensitivity to hyperparameter choices.

\section{Summary of Key Findings}
\label{sec:results_summary}

The experimental results in this chapter address the core research questions posed in the introduction:

\begin{enumerate}
    \item \textbf{Can layout-aware models (LayoutLMv3) reduce hallucination compared to text-only models (Mistral)?}
    \\[2mm]
    \textit{Expected finding}: [To be confirmed] Stage 4 should show dramatically reduced hallucination on sparse layouts (Cluster 0), where Mistral suffers. This validates the hypothesis that spatial grounding prevents spurious predictions.
    
    \item \textbf{Can learned linking improve key-value associations over implicit attention?}
    \\[2mm]
    \textit{Expected finding}: [To be confirmed] Stage 4b (with biaffine linker) should improve over Stage 4a (entity-only), but the magnitude (controlled by $\lambda$ sweep) reveals whether linking is a major bottleneck or a minor sub-problem.
    
    \item \textbf{How much of the improvement is attributable to structural vs.\ architectural changes?}
    \\[2mm]
    \textit{Expected finding}: [To be computed] The Stage 3 vs.\ Stage 4 comparison shows overall improvement, but per-cluster analysis reveals whether improvement is universal or specific to sparse layouts.
\end{enumerate}

Detailed error analysis and per-cluster breakdowns will be presented in appendices (Appendix~\ref{app:error_analysis}).
